\begin{thebibliography}{99}
\bibitem{bright2023}
Paige Bright. \textit{18.S190 Introduction to Metric Spaces}. MIT OpenCourseWare, January IAP 2023. 六讲Lecture Notes及官方TeX, 另附三份Problem Sets. 主讲义PDF物理页数依次为8、8、8、6、7、6.\\
课程: \url{https://ocw.mit.edu/courses/18-s190-introduction-to-metric-spaces-january-iap-2023/}.\\
讲义: \url{https://ocw.mit.edu/courses/18-s190-introduction-to-metric-spaces-january-iap-2023/pages/lecture-notes/}.\\
三份原题: \href{https://ocw.mit.edu/courses/18-s190-introduction-to-metric-spaces-january-iap-2023/pages/problem-sets/}{Problem Sets}, 各3物理页, 包含2题面页与1OCW署名页; 本册第16章逐题译录与编者解答.

\bibitem{rodriguez2021}
Casey Rodriguez(授课), Andrew Lin(笔记). \textit{18.102/18.1021 Introduction to Functional Analysis}, Lecture 3: Quotient Spaces, the Baire Category Theorem and the Uniform Boundedness Theorem. MIT OpenCourseWare, Spring 2021. PDF第3--4页, 正文印刷第16--17页: Baire定理及一致有界原理. 本册重新给出证明, 修正原证明中相对开球与全空间开球的区分.\\
官方资源页: \href{https://ocw.mit.edu/courses/18-102-introduction-to-functional-analysis-spring-2021/resources/lecture-3-quotient-spaces-the-baire-category-theorem-and-the-uniform-boundedness-theorem/}{18.102第3讲}.

\bibitem{munkres2004}
James Munkres(授课). \textit{18.901 Introduction to Topology}, Supplementary Notes D: Countability Axioms. MIT OpenCourseWare, Fall 2004, 3页. 本册可数性讨论的选读参考; 不作为18.901整课或商业教材的全文翻译. Notes G第1--3页为商映射、邻接空间与正规性的参考; Notes C与K保留为选读归档, 未纳入新增核心.\\
\url{https://ocw.mit.edu/courses/18-901-introduction-to-topology-fall-2004/pages/lecture-notes/}.\\
作业: \href{https://ocw.mit.edu/courses/18-901-introduction-to-topology-fall-2004/pages/assignments/}{Assignments}; \href{https://ocw.mit.edu/courses/18-901-introduction-to-topology-fall-2004/d6c5d6d30bf26300b6f1b2c278791796_problemset_5.pdf}{PS5完整原件}, 2物理页(第二页印刷第8页). 本册第18章解答126格; PS0--4及周练只登记官方教材题号, 缺完整题面.

\bibitem{miller2016}
Haynes Miller(授课), Sanath Devalapurkar(课堂LaTeX记录), Xianglong Ni(插图). \textit{18.905 Algebraic Topology I}, MIT OpenCourseWare, Fall 2016. 第5讲物理页1--2: 同伦与形变收缩; 第14讲页1--3: 胞腔粘合; 第18讲页1: 欧拉示性数与同调交错和. 不把后续同调代数与上同调材料称为本册已覆盖.\\
\url{https://ocw.mit.edu/courses/18-905-algebraic-topology-i-fall-2016/pages/lecture-notes/}.

\bibitem{snowden2011}
Andrew Snowden. \textit{18.904 Seminar in Topology}, MIT OpenCourseWare, Spring 2011. Lecture Summaries: 基本群、van Kampen与覆盖的教学安排, 不是完整逐讲证明讲义; 本册相应证明为编者补充.\\
\url{https://ocw.mit.edu/courses/18-904-seminar-in-topology-spring-2011/pages/lecture-summaries/}.

\bibitem{seidel2023}
Paul Seidel. \textit{18.900 Geometry and Topology in the Plane}, MIT OpenCourseWare, Spring 2023. 第4讲印刷页30--33: 绕数; 第29--31讲印刷页219--222、225--228、232--235: 有限复形、边界矩阵与可定向性; 第33讲页246--249: 环路线性计数; 第40讲页292--293: 角亏与离散Gauss--Bonnet. 第40讲为官方公开补充, 当学期课堂未讲. 第30、31、33讲中的符号笔误已在正文修正, 原件保留.\\
\url{https://ocw.mit.edu/courses/18-900-geometry-and-topology-in-the-plane-spring-2023/}.\\
所选CQ原题: \href{https://ocw.mit.edu/courses/18-900-geometry-and-topology-in-the-plane-spring-2023/mit18_900s23_q4.pdf}{第4讲}, \href{https://ocw.mit.edu/courses/18-900-geometry-and-topology-in-the-plane-spring-2023/mit18_900s23_q29.pdf}{第29讲}, \href{https://ocw.mit.edu/courses/18-900-geometry-and-topology-in-the-plane-spring-2023/mit18_900s23_q30.pdf}{第30讲}, \href{https://ocw.mit.edu/courses/18-900-geometry-and-topology-in-the-plane-spring-2023/mit18_900s23_q31.pdf}{第31讲}, \href{https://ocw.mit.edu/courses/18-900-geometry-and-topology-in-the-plane-spring-2023/mit18_900s23_q33.pdf}{第33讲}, \href{https://ocw.mit.edu/courses/18-900-geometry-and-topology-in-the-plane-spring-2023/mit18_900s23_q40.pdf}{第40讲}, 各2物理页(1题面页与1署名页), 共22题. 第31讲第5题另需\href{https://ocw.mit.edu/courses/18-900-geometry-and-topology-in-the-plane-spring-2023/mit18_900s23_lec6.pdf}{第6讲}图(6.9)、(6.11), 原件7物理页, 图在物理第4页、印刷第51页.

\bibitem{mrowka2004}
Tomasz S. Mrowka. \textit{18.965 Geometry of Manifolds}, MIT OpenCourseWare, Fall 2004. 第1、2讲: 图册、光滑映射、基本病态; 第4讲印刷页9--12: 逆函数与隐函数论证, 本册限有限维并补齐条件.\\
\url{https://ocw.mit.edu/courses/18-965-geometry-of-manifolds-fall-2004/pages/lecture-notes/}.

\bibitem{seidel2008}
Paul Seidel. \textit{18.950 Differential Geometry}, MIT OpenCourseWare, Fall 2008. Chapter 4, Lectures 36、39: 长度与内在距离; 本册不译编该章其余测地线与曲率理论. 原件13个物理页, 包括来源页与章扉页.\\
\url{https://ocw.mit.edu/courses/18-950-differential-geometry-fall-2008/pages/lecture-notes/}.

\bibitem{tsinghua2027}
清华大学本科招生网. \textit{2027年丘成桐数学科学领军人才培养计划招生简章}. 2026年9月28日. 第六条第3项: 一试数学1包括分析、代数、几何与拓扑.\\
\url{https://www.admissions.tsinghua.edu.cn/info/1033/2240.htm}.

\bibitem{qzc-syllabus}
清华大学求真书院. \textit{数学领军计划之数学与物理科目考纲发布!}. 官方通知, 2026年8月21日. 本次核实其通知与详细考纲链接存在, 未获得所链原文; 具体考纲核实边界见前言和\texttt{research/syllabus-verification.md}.\\
\url{https://qzc.tsinghua.edu.cn/}. 既有仓库记录的资料题名为\textit{领军考试大纲(2).pdf}, 文件内日期2026年8月20日, 共6页; 本次没有将该摘要提升为重新逐页核验的证据.

\bibitem{cc-license}
Creative Commons. \textit{Attribution--NonCommercial--ShareAlike 4.0 International}.\\
许可: \url{https://creativecommons.org/licenses/by-nc-sa/4.0/}.\\
MIT OCW使用条款: \url{https://ocw.mit.edu/pages/privacy-and-terms-of-use/}.
\end{thebibliography}
